\begin{tikzpicture}[font=\scriptsize, >=latex]
% schematic with measurement points
\draw[thick] (1.0,4.4) -- (1.0,3.4) (1.0,1.6) -- (1.0,0.6);
\draw[fill=white] (0.7,3.4) rectangle (1.3,2.4);
\node at (1.0,2.9) {R1};
\draw[fill=white] (0.7,1.6) rectangle (1.3,0.6);
\node at (1.0,1.1) {R2};
\node[left] at (1.0,4.4) {+5\,V};
\node[left] at (1.0,0.0) {GND};
\fill[accent] (1.0,2.0) circle (0.09);
% probes
\draw[->, thick, accent] (3.4,4.3) -- (1.15,4.35) node[pos=0,below,align=left]{\textbf{P1}: rail really 5 V?};
\draw[->, thick, accent] (3.4,2.2) -- (1.15,2.05) node[pos=0,below,align=left]{\textbf{P2}: the midpoint};
\draw[->, thick, accent] (3.4,0.7) -- (1.15,0.65) node[pos=0,below,align=left]{\textbf{P3}: ground really 0?};
% method box
\draw[thick, fill=goodbg] (4.6,1.4) rectangle (9.4,4.6);
\node[align=left, text width=4.4cm] at (7.0,3.0)
{\textbf{the DMM method}\\[2pt]
1.\ black on a known-good ground\\
2.\ measure the supply first\\
3.\ walk node to node\\
4.\ the first wrong voltage\\
\hspace{8pt}names the dead segment};
\node[align=left, text=slate, text width=9.5cm] at (4.9,-0.7)
{Three probes, one story: P1 proves the rail is live, P3 proves the
return is real, and P2 shows what the divider actually delivered.
Debugging is walking the node list in order and stopping at the first
lie --- never poking at random.};
\end{tikzpicture}
